% titel: Blatt 9, A5: Gerade und ungerade Funktionen
% Teil II: der durchgerechnete Weg.

\subsection*{Gegeben und gesucht}
$f:\mathbb R\to\mathbb R$ stetig; zu zeigen: (a) $f$ gerade
$\Rightarrow\int_{-a}^a f=2\int_0^a f$; (b) $f$ ungerade $\Rightarrow\int_{-a}^a f=0$,
je für alle $a>0$; (c) für eine Stammfunktion $F$ mit $F(0)=0$: $f$ gerade
$\Rightarrow F$ ungerade, $f$ ungerade $\Rightarrow F$ gerade.

\subsection*{a) $f$ gerade}
Mit der Substitution $u=-x$ in $\int_{-a}^{0}$ ($\dd x=-\dd u$, $x=-a\to u=a$,
$x=0\to u=0$):
\begin{align}
  \int_{-a}^{0} f(x)\,\dd x
  = -\int_{a}^{0} f(-u)\,\dd u
  = \int_0^a f(-u)\,\dd u
  = \int_0^a f(u)\,\dd u .
\end{align}
(Letzter Schritt nutzt $f(-u)=f(u)$.) Also
\begin{equation}
  \int_{-a}^{a} f(x)\,\dd x
  = \int_{-a}^{0}f+\int_0^a f
  = \int_0^a f(u)\,\dd u + \int_0^a f(x)\,\dd x
  = 2\int_0^a f(x)\,\dd x .
\end{equation}

\subsection*{b) $f$ ungerade}
Wie in a), aber $f(-u)=-f(u)$:
\begin{equation}
  \int_{-a}^{0}f(x)\,\dd x
  = \int_0^a f(-u)\,\dd u
  = -\int_0^a f(u)\,\dd u
  = -\int_0^a f(x)\,\dd x,
\end{equation}
also $\int_{-a}^a f=\int_0^a f-\int_0^a f=0$.

\subsection*{c) Parität der Stammfunktion}
Da $F$ Stammfunktion mit $F(0)=0$ ist, gilt
$F(x)=\int_0^x f(t)\,\dd t$.

\textbf{$f$ ungerade $\Rightarrow F$ gerade:} mit $u=-t$ ($\dd t=-\dd u$):
\begin{align}
  F(-x)=\int_0^{-x}f(t)\,\dd t
  &= -\int_0^x f(-u)\,\dd u
  = \int_0^x f(u)\,\dd u
   = F(x),
\end{align}
denn $f(-u)=-f(u)$ und $-\cdot(-f(u))=f(u)$. Also $F$ gerade.

\textbf{$f$ gerade $\Rightarrow F$ ungerade:}
\begin{equation}
  F(-x)=-\int_0^x f(-u)\,\dd u
  = -\int_0^x f(u)\,\dd u = -F(x),
\end{equation}
also $F$ ungerade.

\begin{ergebnis}
$f$ gerade: $\displaystyle\int_{-a}^a f=2\int_0^a f$ und $F$ ungerade;
\qquad
$f$ ungerade: $\displaystyle\int_{-a}^a f=0$ und $F$ gerade.
\end{ergebnis}

\begin{kontrolle}
\textbf{Beispiel:} $f(x)=x$ (ungerade) gibt $\int_{-a}^a x\,\dd x=0$, und
$F=\frac{x^2}{2}$ ist gerade. $f(x)=x^2$ (gerade) gibt
$\int_{-a}^a x^2\,\dd x=\frac{2a^3}{3}=2\int_0^a x^2\,\dd x$, und
$F=\frac{x^3}{3}$ ist ungerade. Passt.
\end{kontrolle}

\begin{merksatz}
Die Substitution $u=-x$ macht die Symmetrie rechnerisch sichtbar: gerade
glättet (verdoppelt auf ein halbes Intervall), ungerade hebt auf. Achtung:
Antiderivative \emph{kehrt} die Parität umgekehrt zur Ableitung.
\end{merksatz}
